\section{The just-in-time executor}
\label{sec:executor}

Section~\ref{sec:isa-emulator} states that the execution relation is realised
twice and says what connects the two realisations. This appendix gives the
mechanism of the second one, a just-in-time compiler for Linux x86-64 hosts,
which is what a production prover runs. Nothing in the protocol depends on it:
a platform without the compiler proves the same statements more slowly.

\subsection{Two implementations, three jobs}
\label{sec:exec-roles}

The interpreter is a single loop over \texttt{execute\_cycle}: fetch, dispatch
on the opcode, apply the state transition, append the events the trace
generator will consume. It implements every instruction, every syscall and
every error path, and it is the only implementation that emits the full event
record.

The compiler is used for the two jobs that dominate wall-clock time and need no
events. The first is \emph{whole-program execution}, when a caller wants only
the output stream and the exit status. The second, which the multi-GPU prover
depends on, is \emph{minimal-trace production}: the coordinating process runs
the guest to the end and emits, per shard, a chunk descriptor holding the
pre-access value of every memory word the shard reads, the register records at
its boundary, and the boundary itself. A worker then re-executes the shard from
that descriptor under the interpreter and emits the events locally. The
consequence is that no event record crosses a process boundary and the
coordinator holds one shard's state at a time, and also that the event-emitting
path stays the interpreter's: the compiler never has to reproduce an event
stream, only the state and the accounting a replay needs.

Availability is a build-time decision. The parent crate's build script emits
\texttt{cfg(zkm\_use\_native\_executor)} on Linux x86-64 without the profiling
feature; elsewhere the flag is absent and every entry point stays on the
interpreter, so a platform without the compiler is slower and not otherwise
different.

\subsection{Transpilation and the program cache}
\label{sec:exec-transpile}

Compilation is one-shot and per program, not per basic block discovered at run
time. The driver walks the program's instruction stream once, and for each
instruction records the native offset that the guest program counter maps to,
dispatches on the opcode to a lowering, and emits the epilogue that advances
the clock and the program counter. The result is a native code buffer, mapped
executable from an anonymous file, together with a jump table holding one
native address per guest program counter. Control transfers inside the compiled
program are table dispatches, so the compiled unit is the whole program and a
branch never leaves it.

Two properties of the lowering are worth stating because they are what make it
tractable. Delay slots are resolved statically: a branch or jump whose delay
slot holds another branch or jump, or which ends the program, is rejected at
compile time and the host falls back to the interpreter, so the only
instruction that has to read a delayed target is the one immediately after a
transfer. And an opcode the lowering does not model is a compile-time refusal
rather than a run-time surprise: the driver returns the offending opcode and
the caller falls back.

Compiled programs are cached process-wide. The key must be a hash of the entire
instruction stream: the base address, the instruction count and, for every
instruction, the opcode and all three operands with their immediate flags. A
cheaper fingerprint is unsound in a way that does not announce itself: two
guests that agree on a prefix and a suffix but differ in the middle collide,
and the second one silently executes the first one's code. \Ziren{} hashes the
whole stream for this reason, and a test asserts that changing any single
instruction changes the key.

\subsection{Machine state in host registers}
\label{sec:exec-state}

The guest's register file is held in host registers for the duration of a run.
There are 33 live guest registers (31 general-purpose plus \texttt{HI} and
\texttt{LO}) and 16 vector registers offering 32 half-slots, so exactly one
must live elsewhere: general registers 1 to 30 pack two per vector register
into the first fifteen, \texttt{HI} and \texttt{LO} take the halves of the
sixteenth, and the return-address register takes a general-purpose host
register of its own, because it is read and written by every call sequence and
a half-slot access costs an insert or extract. Register 0 is not stored at all.
The two registers that only \texttt{mmap}-style syscalls touch stay in the
context structure in memory.

Three quantities are pinned to host registers alongside the guest state: the
packed clock, the oracle tail pointer, and the remaining trace-area budget.
All three are written back to the context around any call into the host and on
exit, so the host never reads a stale copy.

Guest memory is a flat array with one sixteen-byte entry per guest word, a
value, the timestamp and shard of its last access, and padding that keeps the
index a shift of the address. The array is a shared mapping of an anonymous
file with reservation disabled, so it is sparse and only touched pages are ever
committed, and an unconstrained block is entered by mapping a private
copy-on-write view of the same file and left by unmapping it, so the rollback
that the paged executor performs with a per-address diff becomes one unmap. Two
alignments in that entry are load-bearing. Its first twelve bytes are exactly
the image of an oracle entry, so recording a pre-access value is one twelve-byte
copy; and its timestamp and shard fields are the little-endian image of the
packed clock register, so stamping an access is one eight-byte store.

\subsection{What the compiled code accounts for}
\label{sec:exec-accounting}

A compiled instruction is not just its state transition. To let a worker replay
a shard, and to close shards where the interpreter would, the emitted code also
performs the bookkeeping of Section~\ref{sec:isa}:

\begin{itemize}
  \item \textbf{Access stamps.} Every register the interpreter stamps is
  stamped here, with the same shard and the same sub-cycle position, at the
  highest position at which the instruction touches it.
  \item \textbf{The clock.} An instruction executes at $\mathsf{clk}$ and its
  accesses occupy $\mathsf{clk}+1,\dots,\mathsf{clk}+4$ by position, so the
  clock advances by five per instruction. The pinned register holds the shard
  and the clock as one word, which is what makes the memory stamp a single
  store.
  \item \textbf{The oracle.} A read or write of an address above the register
  window copies the entry as it stood before the access to the oracle tail,
  which is the descriptor the replaying worker will answer memory reads from.
  \item \textbf{Shard budgets.} The trace-area budget in its pinned register
  and the per-unit height budgets in the context are charged the exact amounts
  the interpreter's accumulator charges, supplied per instruction by the
  compiler. Exhaustion of any one of them is the condition under which the
  interpreter would fence, and the compiled function returns a shard-fence
  status so the host seals the chunk and re-enters.
\end{itemize}

Because the budgets are charged in native code and checked on every
instruction, the shard boundaries a compiled run produces are the boundaries an
interpreted run would produce, which is what allows the two to be compared
chunk by chunk.

\subsection{What the interpreter keeps}
\label{sec:exec-traps}

The compiler deliberately implements no error path and no syscall. Every
\texttt{SYSCALL}, every division by zero, every trap instruction, every
misaligned or out-of-range address, and every operand form the lowering does
not model is a \emph{trap site}: the emitted code stores the guest program
counter and jumps to a shared stub whose host handler synchronises the native
state into the interpreter, runs that one instruction there, synchronises back
and resumes the compiled code at the resulting program counter. The interpreter
therefore owns the precompiles, the unconstrained blocks and every failure, and
the compiler owns the arithmetic, memory and control flow that make up almost
every instruction of a real workload.

This split is what keeps the compiler small enough to argue about. It is also
why a syscall-dense guest sees less speedup than a compute-dense one: each trap
is a state synchronisation in both directions.

\subsection{Evidence that the two agree}
\label{sec:exec-equivalence}

Three checks connect the implementations, and none of them is a proof of
equivalence for every program.

The producer's output is compared against the interpreter's directly: tests
assert that the chunk descriptors a compiled run emits are byte-identical to
those of an interpreted run of the same program, which covers the state, the
oracle, the stamps and the shard boundaries in one comparison. The 770
specification vectors of Section~\ref{sec:verification} are executed on both
implementations and against an independent reference emulator. And continuous
integration compares digests of the emitted records across runs and across
implementations on the tested corpus, which is evidence of determinism and of
agreement there rather than everywhere.

Section~\ref{sec:perf-emu} measures what the compiler buys: on the reference
block it emulates the guest at roughly 130\,MHz against the
interpreter's 5--20\,MHz, depending on whether the interpreter
is recording events.

\subsection{The record pipeline under a host budget}
\label{sec:app-hostbudget}

Section~\ref{sec:perf-emu} states that on a host feeding several accelerators the executor, not the accelerator, is what binds, because each shard is re-emulated by the worker that proves it. The event stream written by one shard replay was approximately 1\,GB. Table~\ref{tab:host} lists the changes made under an explicit host budget. Each preserves the represented execution semantics: the trace columns every constraint reads are unchanged, so the verifying keys and the proofs are unchanged. The record layout itself does change; the table below removes fields and narrows record widths, so each change is witness-equivalent rather than byte-identical, and it is the event record, not the committed trace, whose bytes move.

\begin{table}[t]
\centering
\footnotesize
\setlength{\tabcolsep}{4pt}
\caption{Emulator and record pipeline under a host memory and processor budget. Sizes are per event or per shard of the block of \S\ref{sec:perf-setup}; times are host-side, single process, three repetitions.}
\label{tab:host}
\begin{tabular}{p{0.44\textwidth} r r p{0.27\textwidth}}
\toprule
Change & before & after & effect \\
\midrule
Per-cycle event: keep clock, \textsf{pc} pair and exit code only (no CPU table reads the rest) & 184\,B & 20\,B & event stream per shard 1{,}000\,MB $\to$ 462\,MB; replay 700\,ms $\to$ 435\,ms \\
Memory-instruction event: drop the three fields no column reads & 168\,B & 136\,B & every load/store event \\
Register-access record: carry only the witnessed values (no shard fields; one arm for read/write) & 48 / 24\,B & 20 / 16\,B & every instruction event \\
Memory oracle shipped to a worker: pre-access value, stamp and shard as one byte run & 24\,B, 25\,MB & 12\,B, 12\,MB & per access, per shard; serialise 6--7 $\to$ 1\,ms \\
Coordinator producer without checkpoints: no state clone or first-touch map per shard & 22.9\,s & 20.4\,s & per block; state $O(\text{shard})$ \\
Native (JIT) producer in the coordinator & 12.3\,s & 6.1\,s & per 420\,M-cycle block \\
Flat guest memory: one 16-byte entry per word in a sparse mapping; untouched pages cost nothing; unconstrained blocks are copy-on-write views & paged table & flat, sparse & no per-block memory diff \\
\bottomrule
\end{tabular}
\end{table}

Two rules follow. An event should carry exactly what a column witnesses: three of the seven rows remove fields that only a debug assertion or an unused code path had kept alive; on the shared register records alone this is 44 bytes for every instruction event, one write record and two read records. And the process that coordinates several GPUs should never hold the whole execution: the coordinator's resident state is one shard's oracle and the guest image, so the block size is not bounded by host memory; only the shard size is bounded, by the accelerator's memory (\S\ref{sec:perf-changes}). Both rules apply to any \zkVM{} whose prover is fed by a re-executing host.
